% Unverified replacement proof.  The real-PL sheaf is never identified
% with the real span of global rational model functions in this argument.

\begin{proof}
We use additive weights.  The distance between two continuous metrics is
\[
 d_\infty(\lVert\cdot\rVert_1,\lVert\cdot\rVert_2)
  :=\sup_{C^{\mathrm{an}}}
       \left|\log\frac{\lVert\cdot\rVert_1}
                       {\lVert\cdot\rVert_2}\right|;
\]
the quotient is independent of the frame.

\emph{1. Curvature and degree.}
By \cite[\S1.2 and \S4.1.2]{GR25Classical}, CLD-psh weights are
Thuillier-subharmonic on a smooth strict analytic curve over any complete
nontrivially valued field.  If \(s=f s'\), then
\[
 -\log\lVert s\rVert=-\log\lVert s'\rVert-\log|f|,
\]
and \(\log|f|\) is harmonic.  Thus the weights define one global
subharmonic metric and a positive curvature measure \(\mu\).
The metrized Poincar\'e--Lelong formula and its degree consequence
\cite[Theorem~18.1.6 and Corollary~18.1.7]{CLD25} give
\[
 \mu=\mathrm d'\mathrm d''[-\log\lVert\sigma\rVert]
          +\delta_{\operatorname{div}(\sigma)},\qquad
 \mu(C^{\mathrm{an}})=\deg L
\]
for a nonzero meromorphic section \(\sigma\).  Closed points and component
degrees are counted with their residue-field degrees.  This proves the
first assertion.  If the degree is zero, the positive measure vanishes.
For a local continuous weight \(h\), both \(\mathrm d'\mathrm d''h\)
and \(\mathrm d'\mathrm d''(-h)\) are then zero.  Thuillier's current
criterion \cite[Theorem~3.4.12]{Thu05} says that \(h\) and \(-h\) are
subharmonic.  By \cite[\S3.12]{GR25Classical}, this is exactly
harmonicity, so the metric is harmonic.

\emph{2. Two elementary facts about a finite metric graph.}
Let \(G\) be a finite connected metric graph.  For a continuous function
affine on each edge, use the Laplacian convention
\[
 (\Delta_G f)(\{v\})=\sum_{\nu\in T_vG}\partial_\nu f.
\]
A break in an oriented edge with left slope \(r_-\) and right slope
\(r_+\) contributes \(r_+-r_-\).  Loops contribute both incident
directions.  In this convention the energy identity is
\[
 \int_G |f'|^2=-\int_G f\,\Delta_G f.
\]
It follows by integration by parts separately on the finitely many
edges.  In particular, the graph Laplacian on functions affine between
specified vertices is a symmetric matrix whose kernel consists exactly
of the constants: zero energy forces equality along every edge, and
connectedness then forces constancy.  Its image is the space of
vertex-mass vectors of sum zero.  Hence any finite atomic signed measure
of total mass zero has a piecewise affine potential, unique modulo
constants, after its support has been inserted among the vertices.

We will also use an explicit continuity estimate.  If
\(\Delta_G g=m(\delta_y-\delta_x)\), \(m\geq0\), and \(g(p)=0\) at a
fixed point \(p\), then
\begin{equation}
\label{eq:p081-corrected-transport-bound}
                 \lVert g\rVert_\infty\leq m\,d_G(x,y).
\end{equation}
Indeed, the maximum principle on an edge and at a vertex of zero
Laplacian shows that the maximum and minimum occur at \(x\) and \(y\),
respectively, unless \(g\) is constant.  The energy identity yields
\(\int_G|g'|^2=m\,\operatorname{osc}(g)\).
Along a shortest path from \(x\) to \(y\), the Cauchy--Schwarz inequality
therefore gives
\[
 \operatorname{osc}(g)^2
 \leq d_G(x,y)\int_G|g'|^2
 =m\,d_G(x,y)\operatorname{osc}(g).
\]
This proves the estimate, since the normalization puts zero between
the extreme values.  Potentials transporting finitely many masses
satisfy the sum of the corresponding bounds.

\emph{3. Reduction to a skeletal metric over an algebraically closed
field.}
For the moment suppose that \(K\) is algebraically closed and that
\(d:=\deg L>0\).  By \cite[Theorem~3.4.19]{Thu05}, a continuous
subharmonic metric is the increasing envelope, in norm notation, of
global lisse subharmonic metrics on the same line bundle.  Here
\emph{lisse} is Thuillier's skeletal notion, not ordinary CLD
smoothness.  The proof of that theorem constructs the metrics globally
and checks their transition functions.  The approximation is uniform:
the logarithms of the quotients of the target norm by the approximating
norms are nonnegative continuous functions decreasing to zero on the
compact curve.  For each positive accuracy their sublevel sets form an
increasing open cover.  A finite subcover and directedness give uniform
convergence.  A choice for each accuracy \(2^{-j}\) gives a sequence.
It suffices to approximate one lisse subharmonic metric.

Choose a formal reference metric \(\lVert\cdot\rVert_0\) on \(L\).
Existence and the formal--PL correspondence are
\cite[\S2.2 and Proposition~2.10]{GM19}.  After a common semistable
refinement, the reference curvature \(\mu_0\) is supported on the
type-II vertices of a finite skeleton, and the relative weight of the
lisse metric has the form \(\Phi\circ\tau_G\), with \(\Phi\) piecewise
affine on a finite metric graph \(G\).  The required factorization and
common refinement are \cite[Proposition~3.2.4 and Lemma~2.2.22]{Thu05}.
On the graph its curvature is
\[
                    \mu_\ell=\mu_0+\Delta_G\Phi.
\]
There may initially be type-III breakpoints.  For each such atom
\(m_x\delta_x\) of the positive measure \(\mu_\ell\), move the mass to
a nearby type-II point \(y_x\) on the same edge.  Type-II points are
dense there because \(\Gamma:=-\log|K^\times|\) is nonzero and divisible.
Solve
\[
 \Delta_G g=\sum_xm_x(\delta_{y_x}-\delta_x),\qquad g(p)=0.
\]
By \eqref{eq:p081-corrected-transport-bound}, the norm of \(g\) is at most
\(\sum_xm_xd_G(x,y_x)\), so this change can be made arbitrarily small.
The new metric, with relative weight
\((\Phi+g)\circ\tau_G\), has positive curvature supported only at
type-II points.  An old type-III vertex has valence two and zero
remaining Laplacian, so it is no longer a breakpoint.  Subdivide at
the remaining type-II breakpoints; by \cite[Lemma~2.2.22]{Thu05} this is
a strict semistable refinement over \(K\).

We have therefore reduced to a finite semistable skeleton \(G\), with
vertex set \(V\), edge lengths in \(\Gamma\), and a function \(\Phi\)
affine on every edge, such that
\begin{equation}
\label{eq:p081-corrected-vertex-curvature}
 \mu_0+\Delta_G\Phi=\sum_{v\in V}b_v\delta_v,\qquad
 b_v\geq0,\qquad\sum_{v\in V}b_v=d>0.
\end{equation}
The reference model is included in the refinement.  We have not
identified \(\Phi\) with a real linear combination of global model
functions.

\emph{4. Rational-slope approximation with positive vertex margin.}
Let
\[
        \nu=\frac{d}{|V|}\sum_{v\in V}\delta_v.
\]
The first graph fact gives a potential \(h\) with
\(\Delta_G h=\nu-\sum_v b_v\delta_v\).  For \(0<\lambda<1\), put
\(\Phi_\lambda=\Phi+\lambda h\).  It is still affine on the edges, and
\[
 \mu_0+\Delta_G\Phi_\lambda
   =\sum_{v\in V}\left((1-\lambda)b_v+
                         \lambda\frac{d}{|V|}\right)\delta_v.
\]
Every vertex thus has curvature at least
\(\sigma:=\lambda d/|V|>0\), while
\(\lVert\Phi_\lambda-\Phi\rVert_\infty
 \leq\lambda\lVert h\rVert_\infty\).

We describe explicitly a rational-slope approximation \(\Psi\) of
\(\Phi_\lambda\).  If the graph is a point, approximate its one value by
an element of \(\Gamma\), and there is nothing further to prove.
Otherwise choose
\[
 0<\delta<\frac{\sigma}{2\max(1,\max_{v\in V}|T_vG|)}.
\]
Choose \(q_v\in\Gamma\), close enough to \(\Phi_\lambda(v)\) that, on
each oriented edge \(e\) from \(v\) to \(w\), of length \(\ell_e\), the
number
\[
 s'_e=\frac{q_w-q_v}{\ell_e}
\]
differs from the slope \(s_e\) of \(\Phi_\lambda\) by less than
\(\delta/2\).  This is possible by density of \(\Gamma\) and finiteness
of the graph.  Choose rational numbers
\[
 r_e^-<s'_e<r_e^+,\qquad |r_e^\pm-s'_e|<\delta/2,
\]
and set
\begin{equation}
\label{eq:p081-corrected-breakpoint}
 t_e=\frac{r_e^+\ell_e-(q_w-q_v)}{r_e^+-r_e^-}.
\end{equation}
Then \(0<t_e<\ell_e\).  Divisibility of \(\Gamma\) implies that
\(t_e\in\Gamma\).  On this edge define \(\Psi\) to start at \(q_v\),
have slope \(r_e^-\) until \(t_e\), and have slope \(r_e^+\)
afterwards.  Formula \eqref{eq:p081-corrected-breakpoint} makes its final
value exactly \(q_w\).  The construction also works on a loop, for
which \(q_w=q_v\).  It therefore gives one continuous global function
on the graph, not independently chosen edgewise weights.

The new interior breakpoint has positive Laplacian
\(r_e^+-r_e^-\).  At an original vertex each outgoing slope has changed
by less than \(\delta\).  Thus its curvature is at least
\(\sigma-|T_vG|\delta>\sigma/2\).
It follows that
\[
                         \mu_0+\Delta_G\Psi\geq0.
\]
All vertex and breakpoint values belong to \(\Gamma\), all subdivision
lengths belong to \(\Gamma\), and all slopes are rational.  After
multiplication by one positive integer, \(\Psi\) is integral
\(\Gamma\)-affine.  By \cite[Proposition~3.5(c)]{GJKM19},
\(\Psi\circ\tau_G\) is a \(\mathbf Q\)-model function.  The model slope
criterion \cite[Proposition~3.7]{GJKM19} shows that
\(\lVert\cdot\rVert_0\mathrm e^{-\Psi\circ\tau_G}\) is semipositive.
In that criterion the residue weights are one in the present
algebraically closed setting.

For the approximation bound, put
\(\eta=\max_v|q_v-\Phi_\lambda(v)|\).  The slopes on either piece of
an edge differ from \(s_e\) by less than \(\delta\), so
\[
 \lVert\Psi-\Phi_\lambda\rVert_\infty
       \leq\eta+\delta\max_e\ell_e.
\]
First choose \(\lambda\) small, then \(\delta\) small subject to the
vertex-margin inequality, and finally the \(q_v\) sufficiently close.
The resulting semipositive model metrics approximate the skeletal
metric arbitrarily well.  Together with Step 3 this proves uniform
\(\mathbf Q\)-formal approximation for positive degree over an
algebraically closed field.  The use of \(d>0\) occurs precisely in the
strict vertex margin; it cannot be suppressed at degree zero.

\emph{5. Descent of the positive-degree approximation to \(K\).}
Return to an arbitrary \(K\), still with \(\deg L>0\), and put
\(F=\widehat{K^{\mathrm a}}\).  Smoothness and geometric connectedness
imply that \(C\) is geometrically integral.  Its base change to \(F\)
is a smooth proper connected curve.  Scalar extension preserves
CLD-psh weights, including their transfinite stages, by
\cite[Proposition~8.1.5 and \S11.1.9]{CLD25}.  Step 4 therefore supplies
semipositive \(\mathbf Q\)-formal approximants over \(F\).

Choose a formal reference metric over \(K\), and write
\[
 \theta=c_1(L,\lVert\cdot\rVert_0),\qquad
 u=-\log\frac{\lVert\cdot\rVert}{\lVert\cdot\rVert_0}.
\]
The metric--function dictionary is \cite[\S4.7]{GM19}.  Choose upstairs
\(\theta_F\)-psh model functions \(\varphi_i^F\) with
\(\|\varphi_i^F-u_F\|_\infty<i^{-1}\).
By \cite[Proposition~2.18(b)]{GM19}, each upstairs
\(\mathbf Q\)-PL metric is defined over a finite subextension
\(K_i/K\) of \(F/K\).  It is semipositive over \(K_i\): after clearing
denominators choose a representing formal model using
\cite[Proposition~2.10]{GM19}; its base change is nef by
\cite[Proposition~3.5]{GM19}, and nefness descends by
\cite[Lemma~3.3]{GM19}.  Denote the descended model function by
\(\varphi_i^{K_i}\).

Choose
\[
 c_i\in\mathbf Q(-\log|K_i^\times|),\qquad
                     i^{-1}\leq c_i\leq2i^{-1},
\]
and put \(\psi_i^{K_i}=\varphi_i^{K_i}-c_i\).  Such constants exist
because a nonzero divisible subgroup of \(\mathbf R\) is dense.
They are \(\mathbf Q\)-model constants, so the shift does not change
semipositivity.  Upstairs,
\[
                         u_F-3i^{-1}\leq\psi_i^F\leq u_F.
\]
Take a finite normal extension \(N_i/K\) in \(K^{\mathrm a}\)
containing \(K_i\), for example the splitting field of finitely many
generating minimal polynomials.  Separability is not required.
Write \(q_i:C_{N_i}^{\mathrm{an}}\to C^{\mathrm{an}}\).
Surjectivity of scalar-extension projections
\cite[\S11.1.9]{CLD25} gives
\[
 q_i^*u-3i^{-1}\leq\psi_i^{N_i}\leq q_i^*u.
\]
Consequently
\[
 q_i^*u-3i^{-1}
 \leq P_{q_i^*\theta}(q_i^*u)\leq q_i^*u.
\]
For the finite normal extension, and only for that extension,
\cite[Proposition~2.11]{GJKM19} identifies the middle expression with
\(q_i^*P_\theta(u)\).  Surjectivity of \(q_i\) and passage to the limit
give
\begin{equation}
\label{eq:p081-corrected-envelope}
                              P_\theta(u)=u.
\end{equation}
Here the class \(c_1(L)\) is ample, since \(\deg L>0\).
The ample curve envelope theorem
\cite[Theorem~1.1 and Proposition~2.9(viii)]{GJKM19}, together with
\eqref{eq:p081-corrected-envelope}, gives uniform approximation by
\(\theta\)-psh model functions over \(K\).  Their metrics are the
required semipositive \(\mathbf Q\)-formal metrics on the fixed
line bundle \(L\).  This use of the envelope theorem is exclusively in
positive degree; no nonemptiness assertion is inferred from a
degree-zero numerical class.

\emph{6. Relatively ample models and scalar extensions.}
By \cite[Theorem~1.1]{GM19}, semipositive formal and algebraic model
metrics agree on a proper algebraic variety.  Represent one approximant
by a nef \(\mathbf Q\)-model \(\mathcal N_j\).  Proposition~4.11 of
\cite{GM19} provides a dominating algebraic model on which a positive
tensor power of \(L\) has an ample extension.  Divide by that tensor
power and call the resulting ample \(\mathbf Q\)-model \(\mathcal A_j\).
The pullback of \(\mathcal N_j\) is nef.

For a sufficiently small rational \(0<\varepsilon_j<1\), the
\(\mathbf Q\)-model
\[
            (1-\varepsilon_j)\mathcal N_j+
                           \varepsilon_j\mathcal A_j
\]
has metric at distance less than \(j^{-1}\) from the original
approximant.  Its degree on each special-fibre component is strictly
positive.  Hence it is relatively ample by the fibrewise criterion
used in \cite[Proposition~4.13]{GM19}.  These relatively ample metrics
still converge to the target.  Clear denominators and take a tensor
power that is relatively generated.  If \(s_0,\ldots,s_N\) are model
sections generating that power, its model metric is characterized by
\[
                           \max_\alpha\|s_\alpha(x)\|=1.
\]
This is the Fubini--Study metric of the generated \(K^\circ\)-lattice.
Taking the rational root proves the Fubini--Study assertion.

Under any complete extension \(E/K\), model metrics remain model
metrics and nefness is preserved
\cite[Proposition~2.18(a) and Lemma~3.3]{GM19}; relative ampleness is
also preserved by base change.  Surjectivity of
\(C_E^{\mathrm{an}}\to C^{\mathrm{an}}\) shows that pullback preserves
the uniform metric distance.  The pulled-back sequence therefore
proves all the extension assertions in
\cref{itm:p081-corrected-positive}.

\emph{7. The exact degree-zero criterion.}
Suppose that \(\deg L=0\).
If a sequence of semipositive \(\mathbf Q\)-formal metrics exists,
then any member supplies the model metric required in the statement.
Conversely, let \(\lVert\cdot\rVert_*\) be one such metric.
The model positivity criterion and its curve subharmonicity
interpretation \cite[Proposition~3.5]{GM19},
\cite[Proposition~5.1]{GR25PL} give a positive curvature measure.
Its mass is zero, so it too is harmonic.  Hence
\[
                  a=-\log
                    \frac{\lVert\cdot\rVert}{\lVert\cdot\rVert_*}
\]
is a global continuous harmonic function on the proper connected curve.
It is constant: it attains a maximum by compactness, and the harmonic
maximum principle makes its maximum locus open as well as closed.
Choose constants
\(a_j\in\mathbf Q(-\log|K^\times|)\) tending to \(a\).
The metrics
\(\lVert\cdot\rVert_*\mathrm e^{-a_j}\) are semipositive
\(\mathbf Q\)-formal metrics converging uniformly to the given one.

They are numerically trivial on the special fibre.  To see this without
assuming a discrete valuation or reduced special fibre, choose an
algebraic nef model of a positive tensor power, say \(\mathcal L\)
extending \(L^{\otimes m}\).  For \(0\ne\rho\in K^{\circ\circ}\), the
rank-one divisorial degree formula of \cite[\S4.8]{GM19} gives
\[
 (-\log|\rho|)m\deg L
      =\sum_Y\lambda_Y(\rho)\deg_Y(\mathcal L),\qquad
                            \lambda_Y(\rho)>0.
\]
Here \(Y\) runs over the reduced irreducible special-fibre components;
the coefficients come from
\(\operatorname{cyc}(\operatorname{div}(\rho))\), not an assumption
that the special fibre is itself a Cartier divisor.  The left side
vanishes.  Nefness makes each component degree nonnegative, so each is
zero.  Every proper integral curve in the special fibre is one of
these components, proving numerical triviality.  The constant shifts
preserve these degrees.  The same scalar-extension argument as in
Step 6 completes \cref{itm:p081-corrected-zero}.
\end{proof}
